\newpage
\phantomsection
{\centering \section*{KẾT LUẬN}}
\addcontentsline{toc}{section}{KẾT LUẬN}

\subsection*{Kết luận}

Sau thời gian nghiên cứu,
thiết kế và phát triển,
đồ án tốt nghiệp
đã hoàn thành đầy đủ các mục tiêu đề ra ban đầu,
đáp ứng toàn bộ các yêu cầu kỹ thuật
và nghiệp vụ được giao.

Về kết quả đạt được, đồ án đã mang lại các giá trị cụ thể như sau:

\begin{itemize}
\item \textbf{Về mặt sản phẩm:}
Xây dựng thành công một hệ thống trợ lý pháp luật hoàn chỉnh
hoạt động dựa trên kiến trúc vi dịch vụ
ổn định và dễ mở rộng.
Trợ lý AI sử dụng giải pháp Agentic RAG tiên tiến với LangGraph,
cho phép trả lời các truy vấn pháp lý tiếng Việt một cách chính xác,
minh bạch qua luồng suy luận rõ ràng và dẫn nguồn từ CSDL pháp điển cũng như văn bản quy phạm pháp luật chính thống.
Đồng thời, hệ thống hỗ trợ đàm thoại giọng nói thời gian thực, xử lý tài liệu người dùng và tự động hóa quy trình thanh toán.

\item \textbf{Về mặt công nghệ:}
Tích hợp và làm chủ nhiều giải pháp công nghệ hiện đại:
xây dựng hệ thống trên nền tảng kiến trúc vi dịch vụ
có ứng dụng Domain-Driven Design,
vận hành qua Kong API Gateway
với các giao thức giao tiếp REST, gRPC
cùng hệ thống hàng đợi RabbitMQ, Kafka.
Thiết lập lưu trữ đa dạng sử dụng postgres
kết hợp các CSDL vector (Qdrant, Pinecone, Milvus)
cùng các kho lưu trữ Supabase Storage, Cloudflare R2, Google Drive.
Xây dựng các đường ống thu thập dữ liệu
tự động nhờ cron-job.org, GitHub Actions, Scrapy và dlt.
Triển khai giải pháp Agentic RAG nâng cao bằng LangGraph
kết hợp các mô hình ngôn ngữ lớn
và HuggingFace (vietnamese-sbert).
Tích hợp hệ thống xác thực bảo mật qua Supabase (JWT, Google OAuth, 2FA/TOTP)
cùng phân hệ đàm thoại giọng nói LiveKit
và các công cụ TTS (edge-tts, piper-tts, ElevenLabs).
Toàn bộ hệ thống được đóng gói bằng Docker,
điều phối vận hành trên Kubernetes
thông qua quy trình CI/CD của GitHub Actions
và đồng bộ hóa hạ tầng tự động theo mô hình GitOps với ArgoCD.

\item \textbf{Về mặt kiến thức và kỹ năng tích lũy:}
Qua quá trình thực hiện đồ án,
sinh viên đã tích lũy được nhiều kiến thức
về
mô hình ngôn ngữ lớn,
kiến trúc hệ thống phân tán
và quy trình vận hành DevOps/GitOps.
Đồng thời, sinh viên phát triển
và nâng cao các kỹ năng như:
kỹ năng nghiên cứu khoa học,
tìm kiếm và tổng hợp tài liệu,
kỹ năng lập kế hoạch và quản lý dự án.

\end{itemize}

Về mặt đánh giá,
hệ thống cho thấy tính thực tiễn
khi giải quyết được bài toán tra cứu
thông tin pháp luật phức tạp.
Tuy nhiên, tính chính xác của câu trả lời phụ thuộc nhiều
vào chất lượng và độ cập nhật của CSDL nguồn.
Các mô hình AI vẫn đóng vai trò
là một trợ lý hỗ trợ đắc lực
và không thể thay thế hoàn toàn vai trò của các chuyên gia pháp lý
trong các tình huống quyết định mang tính pháp lý cao.

\subsection*{Hướng phát triển}

Để hoàn thiện ứng dụng
và nâng cao khả năng
phục vụ người dùng trong thực tế,
hướng phát triển của đồ án
trong tương lai tập trung vào các nội dung sau:

\begin{itemize}

\item \textbf{Nâng cao chất lượng AI và tối ưu hóa hệ thống RAG:}
Nghiên cứu
huấn luyện tinh chỉnh
mô hình ngôn ngữ lớn
trên các tập dữ liệu đặc thù của pháp luật Việt Nam
để nâng cao độ chính xác
và khả năng lập luận chuyên sâu.

\item \textbf{Nâng cấp khả năng điều phối luồng dữ liệu:}
Tích hợp Airflow vào quy trình xử lý dữ liệu
để tự động hóa và giám sát đường ống dữ liệu.
Sử dụng đồ thị định hướng không chu trình
(Directed Acyclic Graph - DAG)
để định nghĩa các luồng logic phức tạp dựa trên điều kiện thực tế của dữ liệu.

\item \textbf{Đánh giá và kiểm soát chất lượng mô hình AI:}
Ứng dụng khung đánh giá RAGAS
(Retrieval Augmented Generation Assessment)
để đo lường các chỉ số chất lượng của hệ thống RAG.
Từ đó, tiến hành tinh chỉnh và cấu hình các ngưỡng phản hồi
hợp lý nhằm đảm bảo tính an toàn,
chính xác và đáng tin cậy của mô hình
trước khi trả về kết quả cho người dùng.

\item \textbf{Đảm bảo chất lượng phần mềm:}
Triển khai công cụ Playwright
để xây dựng các kịch bản kiểm thử tự động
toàn diện cho cả giao diện người dùng
và các API,
giúp phát hiện lỗi sớm
và đảm bảo hệ thống hoạt động ổn định.

\item \textbf{Mở rộng phân hệ thanh toán và tài chính:}
Phát triển thêm tính năng xuất hóa đơn
điện tử phục vụ đối tượng khách hàng doanh nghiệp
và xây dựng chức năng mã giảm giá cần
xử lý thêm bài toán nhiều người
sử dụng mã giảm giá cùng lúc,
đồng thời tích hợp thêm cổng thanh toán Stripe
nhằm đáp ứng nhu cầu của tệp khách hàng quốc tế.

\item \textbf{Tối ưu hóa trải nghiệm đàm thoại giọng nói (Voice Bot):}
Nâng cấp các dịch vụ chuyển đổi giọng nói
cao cấp
hoặc tự triển khai các mô hình AI trên máy chủ có hỗ trợ GPU
nhằm giảm thiểu tối đa độ trễ, mang lại trải nghiệm đàm thoại mượt mà và tự nhiên cho người dùng.

\item \textbf{Nâng cao khả năng xử lý tài liệu đa định dạng:}
Tích hợp các công nghệ trích xuất thông tin
và đọc hiểu tài liệu tiên tiến
để hệ thống có thể xử lý đa dạng các định dạng tệp tin.

\item \textbf{Mở rộng và tối ưu hóa hạ tầng điện toán đám mây:}
Triển khai
trên môi trường đa cụm Kubernetes
nhằm tăng cường tính sẵn sàng,
đồng thời tối ưu hóa các quy trình tự động co giãn
giúp hệ thống vận hành ổn định trước lượng lớn người dùng truy cập đồng thời.

\item \textbf{Xây dựng hệ thống giám sát và quản lý nhật ký tập trung (Observability):}
Triển khai giải pháp giám sát toàn diện cho kiến trúc vi dịch vụ
bằng cách tích hợp Prometheus để thu thập số liệu hiệu năng (metrics)
hệ thống theo thời gian thực
và Grafana để trực quan hóa trạng thái vận hành, thiết lập các kịch bản cảnh báo tự động.
Đồng thời, ứng dụng bộ công cụ ELK Stack
(Elasticsearch, Logstash, Kibana)
làm giải pháp quản lý nhật ký (logging) tập trung,
hỗ trợ thu thập, tiền xử lý và phân tích log
từ tất cả các vi dịch vụ, từ đó tối ưu hóa quy trình
định vị sự cố (troubleshooting) và đảm bảo tính sẵn sàng cao cho hệ thống.

\end{itemize}
\newpage
